\section{研究目标与系统设计}
\label{sec:system}
\subsection{把一次程序优化组织为连续研究}
一次程序优化从可运行的初始版本开始。研究者提出改动，用相同任务检验正确性和目标指标，再决定保留、撤销或继续修改。生成模型能够提供多种改法，而研究能推进多远，还取决于它把多少请求用于分析、何时需要更强的模型，以及何时结束搜索。QuotaCrux（枢决）将这些选择放入同一个可执行循环，并记录每次选择的结果与资源消耗。

系统接收任务契约、初始程序 $c_0$、评测接口和资源预算。研究过程可以写为状态序列
\begin{equation}
 s_t=(\mathcal{T},\mathcal{C}_t,\mathcal{E}_t,\mathcal{H}_t,\mathbf{b}_t),
\end{equation}
其中 $\mathcal{T}$ 为任务，$\mathcal{C}_t$ 为候选集合，$\mathcal{E}_t$ 为实验结果，$\mathcal{H}_t$ 为行动历史，$\mathbf{b}_t$ 为剩余预算。每一步选择行动、执行并获得新证据，随后构造下一状态。这里的改进发生在同一任务的候选程序和搜索策略上。本轮实验均从空经验快照 \texttt{v0} 启动，以隔离单轮研究内的搜索与调度效应（跨任务经验复用与历史复验路径不在本轮评测范围）。

\subsection{模块职责与外部判断}
Runner 负责连接任务、模型和工具：构造可见状态，执行结构化行动，检查预算，保存候选并推进研究。任务包固定输入接口、初始代码、目标函数和开发／最终数据身份；基准适配器（benchmark adapter）将不同任务接入统一评测通道。ALE 返回合法性与原始分数，AlgoTune 返回数值正确性与函数计时。模型提出的目的、假设和预期结果随行动记录，裁判结果则独立写入运行账本。

\par\noindent
\begin{tabularx}{\linewidth}{@{}lX@{}}\toprule
模块 & 职责与交付内容\\\midrule
任务与 adapter & 固定输入输出契约、数据身份、裁判和目标指标\\
Runner & 执行动作、控制预算、记录状态与停止原因\\
动作／模型路由 & 分别选择下一步行动和承担生成的模型角色\\
生成模型 & 输出分析或完整候选程序，声明预期改进\\
裁判与选择器 & 实测合法性和目标值，保留有效的较好候选\\
记录与导出 & 关联步骤、候选哈希、模型调用、评测和费用\\\bottomrule
\end{tabularx}\par

在本轮原生任务中，行动包括分析、提出候选、重新验证和结束。提出候选后，程序立即进行静态检查与开发评测；重新验证则对已存在的候选再次取得证据。行动记录保留候选父版本，使一次修改能够追溯到它所依据的程序。研究结束时，选择器冻结开发集上的最佳产物，随后执行独立最终测试并导出结果。

\subsection{动作选择与模型选择的解耦}
动作路由回答“接下来做什么”，模型路由回答“这一步交给谁生成”。默认的 LLM 动作路由由当前生成模型同时选择行动和返回内容。Decisions 动作路由先选择离散行动：分析、提出候选、验证某个已有候选或结束；选择分析或提案后，再请求生成模型完成相应内容。这样，动作分配可以与代码生成分别观察。

Decisions 使用 \texttt{gpt-6-luna}，接收当前可见状态、目标方向、剩余预算、候选与生成模型配置。动作问题要求它依据进展选择下一步；模型问题要求它权衡任务难度、已有尝试、剩余时间，以及使用较大模型的费用和延迟。接口返回离散选项、置信度和各选项概率；Runner 执行选项对应的分支，并将概率与调用标识一并归档。

生成角色从 small 开始，允许单向升级到 large。本轮原生规则路由以累计两次候选尝试或三次 small 请求为升级条件，计入的候选尝试包括已经获得改进的候选。Decisions 模型路由则在待生成请求前选择继续 small 或升级 large；升级后，本次研究持续使用 large。固定模型对照始终由 Qwen 生成，路由实验分别替换这两个选择点。

路由也消耗预算。双 Decisions 完成一次提案时，可能依次使用动作选择、模型选择和代码生成三个请求。Runner 根据下一步所需请求数构造菜单，配额不足时移除生成选项；模型选择前为路由和生成预留额度，生成前再次检查剩余时间。由此，决策的收益与决策本身的开销共同体现在一次研究的账本中。

\subsection{三种研究策略改变哪些信息}
F、I、N 使用相同任务和外部选择器，通过状态构造改变模型能够利用的信息。
\par\noindent
\begin{tabularx}{\linewidth}{@{}lXXX@{}}\toprule
可见信息／行动 & F：反馈研究 & I：独立探索 & N：隐藏后续反馈\\\midrule
初始程序及其评测 & 可见 & 可见 & 可见\\
前序分析与候选 & 可见 & 每次从 $c_0$ 开始 & 保留自身历史\\
后续裁判反馈 & 可见 & 隐藏 & 隐藏\\
可选行动 & 分析、提案、复验、结束 & 提案、结束 & 分析、提案、复验、结束\\\bottomrule
\end{tabularx}\par

F 可以根据正确性、分数和失败类别改进已有方案。I 每次独立从初始程序提出完整候选，由外部选择器在全部尝试中选优。N 保留自身代码与分析，隐藏后续评测结果、排序和实际最佳候选；协议格式错误与剩余预算继续可见。F 与 I 比较两种完整搜索方式，F 与 N 则考察保留自身历史时加入实验反馈所产生的变化。
\source{[S1] Runner、prompt、routing 与 Decisions 客户端；原生路由冻结配置。}
